\section{凸性的基本不等式}

设 X 是一个集；在向量空间 \(R^{X}\) （由定义在 X 上的全体有限数值函数 \(^{1}\) 组成）中，设 P 是 X 上全体正实值函数所成的集。另一方面，设 M 是定义在 P 上的一个数值函数 \(^{2}\) ，有限与否皆可，取值 \(\geqslant0\) ，并且满足：

\(1^{\circ} M(0) = 0\) ，且 \(M\) 是正齐次的，也就是说，\(M(\lambda f) = \lambda M(f)\) 对每个实数 \(\lambda > 0\) 成立。

\(2^{\circ} M\) 在 P 中递增，换言之，关系 \(f \leqslant g\) 蕴涵 \(M(f) \leqslant M(g)\) 。

\(3^{\circ} M\) 在 P 中是凸的，换言之，它 (TVS, II, §2, No. 8) 满足关系 \(M(f + g) \leqslant M(f) + M(g)\) 。

例. --- 设 X 是一个有限集，例如 N 的区间 \([1,n]\) ；用 \(x_{i}\) ( \(1 \leqslant i \leqslant n\) ) 表示向量 \(x \in R^{n}\) 的坐标，则函数 \(M_{1}(\mathbf{x}) = \sum_{i=1}^{n} x_{i}\) 与 \(M_{\infty}(\mathbf{x}) = \sup_{1 \leqslant i \leqslant n} x_{i}\) 在由坐标 \(\geqslant 0\) 的向量 x 所组成的集 P 中满足前述条件。

注. --- 设 S 是含于 P 中的一个带顶点的凸锥（也就是说，一个满足 \(S + S \subset S\) 与 \(\lambda S \subset S\) （对 \(\lambda \geqslant 0\) ）的集；参见 TVS, II, §2, No. 4）；设 M 是定义在 S 上、取值 \(\geqslant 0\) 的数值函数（有限与否皆可），且在 S 中满足上述条件 \(1^{\circ}\) 、\(2^{\circ}\) 与 \(3^{\circ}\) 。那么 M 可以延拓到整个集 P 上，并使延拓后的函数（我们仍记作 M）满足同样的条件：对每个函数 \(f \in P\) ，令 \(M(f) = +\infty\) ，如果不存在任何函数 \(g \in S\) 使得 \(f \leqslant g\) ；而在相反的情形，令 \(M(f) = \inf_{g \in S, f \leqslant g} M(g)\) ；如此即

可完成延拓。这一程序将在第 IV 章 §1 中用来定义正函数的上积分。

命题 1. --- 设 \(\varphi(t_{1}, t_{2}, \ldots, t_{n})\) 是一个有限数值函数，对 \(t_{i} \geqslant 0\) ( \(1 \leqslant i \leqslant n\) ) 有定义且连续，并满足：

\(1^{\circ}\) 关系 \(t_i > 0\) \((1\leqslant i\leqslant n)\) 蕴涵 \(\varphi (t_1,t_2,\dots ,t_n) > 0\) ；

\(2^{\circ}\) 函数 \(\varphi\) 是正齐次的；

\(3^{\circ}\) 集 \(\mathbf{K} \subset \mathbf{R}^{n}\) （由关系 \(t_{i} \geqslant 0\) ( \(1 \leqslant i \leqslant n\) ) 及 \(\varphi(t_{1}, t_{2}, \ldots, t_{n}) \geqslant 1\) 定义）是凸的。

在这些条件下，若 \(f_1, f_2, \ldots, f_n\) 是 \(n\) 个有限函数 \(\geqslant 0\) ，定义在 \(X\) 上，且使得 \(M(f_i) < +\infty\) 对 \(1 \leqslant i \leqslant n\) 成立，则

\[
M \left(\varphi \left(f _ {1}, f _ {2}, \dots , f _ {n}\right)\right) \leqslant \varphi \left(M \left(f _ {1}\right), M \left(f _ {2}\right), \dots , M \left(f _ {n}\right)\right).\tag{1}
\]

由 Hahn-Banach 定理 (TVS, II, §5) 可知，K 是 n 个半空间 \(t_{i} \geqslant 0\) ( \(1 \leqslant i \leqslant n\) ) 与一族闭半空间 \((\mathrm{U}_{\iota})_{\iota \in \mathrm{I}}\) 的交，其中 \(U_{\iota}\) 由形如下式的关系定义

\[
\alpha_ {\iota 1} t _ {1} + \alpha_ {\iota 2} t _ {2} + \dots + \alpha_ {\iota n} t _ {n} - \beta_ {\iota} \geqslant 0,\tag{2}
\]

其中诸 \(\alpha_{\iota k}\) 不全为零。根据假设，若 \(\mathbf{t}=(t_{i})\) 满足 \(t_{i}>0\) （对 \(1\leqslant i\leqslant n\) ），则 \(\varphi(t_{1},\ldots,t_{n})>0\) ，因此存在 \(\lambda_{0}>0\) ，使得关系 \(\lambda\geqslant\lambda_{0}\) 蕴涵 \(\lambda t\in K\) ；这表明，对每个 \(\iota\in I\) ，关系 \(t_{i}\geqslant0\) ( \(1\leqslant i\leqslant n\) ) 蕴涵 \(\alpha_{\iota1}t_{1}+\cdots+\alpha_{\iota n}t_{n}\geqslant0\) ，从而 \(\alpha_{\iota k}\geqslant0\) 对 \(1\leqslant k\leqslant n\) 成立；于是显然，K 也是诸半空间 \(t_{i}\geqslant0\) ( \(1\leqslant i\leqslant n\) ) 与诸 \(U_{\iota}\) （即使 \(\beta_{\iota}\geqslant0\) 的那些）的交；此外，由于原点不属于 K，至少存在一个指标 \(\iota\) 使得 \(\beta_{\iota}>0\) 。

现在设 C 是 \(R^{n+1}\) 中由关系 \(t_{i} \geqslant 0\) ( \(1 \leqslant i \leqslant n + 1\) ) 与 \(t_{n+1} \leqslant \varphi(t_{1}, t_{2}, \ldots, t_{n})\) 所定义的凸锥（即在 \(R^{n+1}\) 中由凸集 \(K \times \{1\}\) 生成的凸锥的闭包）；立即可以看出，C 也可以由关系 \(t_{i} \geqslant 0\) ( \(1 \leqslant i \leqslant n + 1\) ) 以及

\[
\beta_ {\iota} t _ {n + 1} \leqslant \alpha_ {\iota 1} t _ {1} + \dots + \alpha_ {\iota n} t _ {n} \quad (\iota \in \mathrm{I}, \beta_ {\iota} \geqslant 0).\tag{3}
\]

因此，对每个 \(x \in X\) ，我们有

\[
\beta_ {\iota} \varphi (f _ {1} (x), \dots , f _ {n} (x)) \leqslant \alpha_ {\iota 1} f _ {1} (x) + \dots + \alpha_ {\iota n} f _ {n} (x)\tag{4}
\]

对一切 \(\iota \in I\) 成立。对每个指标 \(\iota\) （它满足 \(\beta_{\iota} > 0\) ），由 (4) 式及关于 \(M\) 的假设可知 \(M(\varphi(f_1, f_2, \ldots, f_n))\) 有限，且

\[
\beta_ {\iota} M \left(\varphi \left(f _ {1}, f _ {2}, \dots , f _ {n}\right)\right) \leqslant \alpha_ {\iota 1} M \left(f _ {1}\right) + \alpha_ {\iota 2} M \left(f _ {2}\right) + \dots + \alpha_ {\iota n} M \left(f _ {n}\right),
\]

而当 \(\beta_{\iota} = 0\) 时，这一关系也以显然的方式成立。这样一

来即可看出，坐标为

\[
M (f _ {1}), M (f _ {2}), \dots , M (f _ {n}), M (\varphi (f _ {1}, f _ {2}, \dots , f _ {n}))
\]

的点属于 C，这就证明了本命题。

